\section{Angular determinant and Weinberg electroweak relation}

The determinant of \(T\) is

\begin{equation}
\det T=q_+q_-=D_A.
\label{eq:detT}
\end{equation}

The constitutive sector of the vacuum evaluates \(F(T)^2\), where
\(F(z)=(1-z^{90})/(1-z^{120})\), whereas the electroweak angular
sector evaluates \(\det T\). These are two functions of a
common antecedent, not a horizontal factorization between constants.
This distinction is developed in the \cref{ch:ew}.

\begin{remark}[Negative control]
The block is refuted if any of these exact identities is violated:
positivity of \(I-T^{120}\), equality
\cref{eq:three-four-completion}, relations \cref{eq:vacuum-relations}, or
compatibility \(E_m(\mathcal V_{m+1})=\mathcal V_m\). An SI coincidence
cannot compensate for an operator failure.
\end{remark}

\subsection*{Logical Scope and Domain of Validity}
Operator, spectrum, completion \(3\to4\), parity, amplification, and constitutive sections: \emph{Exact internal theorem}. The formulation that unites them in a single operator is \emph{New formalization}; the angular data and lines are an \emph{Antecedent demonstrated result} of the preceding HMT development and are declared here as explicit initial structure.

